\section{Incidence structure and marginal restrictions}\label{sec:incidence}
The \emph{incidence graph} of a specified factorization has variable vertices
$i$, factor vertices $e$, and one edge $ie$ for each $i\in e$. It is a
simple bipartite graph even when two factors have identical supports. This
is different from the primal variable graph, which joins two variables
occurring together in a factor. All widths below concern the incidence graph.
Its treewidth is the minimum, over tree decompositions, of the largest bag
size minus one; a tree decomposition covers every edge and has connected
bags for each vertex. Its degeneracy is the largest minimum degree of a
nonempty subgraph. Its orientation parameter is the minimum possible maximum
outdegree, with no acyclicity requirement. Its variable frequency $\nu$ is
the largest number of factors containing one variable. A feedback variable
set is a set of variable vertices whose deletion leaves a forest; write
$f_F$ for its size. Affine factors are discarded before these parameters
are measured, unless a factorization is explicitly retained.

\subsection{Ownership and asymmetric orientations}
Write $\kappa=2/(1-e^{-1})$. Let $B(d)$ be any valid universal degree-$d$
upper bound from Section~\ref{sec:positive-universal}, with $B(1)=0$.
In particular one can choose $B(d)\sim\log d/\log\log d$.

\begin{theorem}[Incidence orientation bound]\label{thm:incidence-orientation}
For a positive multilinear polynomial on a unit cube, suppose an incidence
orientation has variable outdegrees at most $r$ and factor outdegrees at
most $s$, where $r,s$ are nonnegative integers. Then
\[
 \tgap\le\bigl[r+\min\{s,B(s+1)\}+\kappa\bigr]\hgap.
\]
In particular frequency $\nu$ gives $\nu+\kappa$, and an orientation with
all outdegrees at most $k$ gives $k+\min\{k,B(k+1)\}+\kappa$.
These conclusions transfer to boxes with zero lower bounds by scaling.
\end{theorem}
\begin{proof}
Call a coordinate low if its mean is at most $1/2$. A term with exactly
one low coordinate will be called a one-low term. All other terms receive
at least $t_e/\kappa$ deficiency from independence by
Lemma~\ref{lem:easy-terms}. In a one-low term the low coordinate is its
anchor, of mean $u_e$. Write $q_i=1-x_i$ for the high-coordinate failure
means. Partition its high variables into $I_e$, whose incidence edges
point into the factor, and $O_e$, whose edges point out. Put
\[
 t_e^I=\min\{u_e,\sum_{i\in I_e}q_i\},\qquad
 t_e^O=\min\{u_e,\sum_{i\in O_e}q_i\}.
\]
Capped sums are subadditive, so $t_e\le t_e^I+t_e^O$.

For $r>0$, independently at each high variable label its incidences into
one-low factors by distinct labels in $\{1,\ldots,r\}$. In round $h$ a
variable is owned by its factor with label $h$, if any. Thus the owned
sets $J_{e,h}$ are disjoint over factors and partition $I_e$ over rounds.
Take one uniform $U$ and give every low coordinate its success interval
$[0,x_i]$. For the high variables owned by $e$, arrange failure sets so
that their union inside $[0,u_e]$ has length
$\min\{u_e,\sum_{i\in J_{e,h}}q_i\}$. This can be done by consecutive
arcs of lengths $\min(q_i,u_e)$ on the circle of circumference $u_e$;
when $q_i>u_e$, add measure $q_i-u_e$ outside that interval. For $u_e=0$
use arbitrary failure sets of the required measures. Unowned variables
also receive arbitrary marginal-correct failure sets. Disjoint ownership
makes these assignments compatible, including when several terms share
one low anchor. Their full deficiencies dominate the owned failure events.
A uniform mixture of rounds therefore gives
\[
 \E D_e\ge\frac1r\sum_h\min\{u_e,\sum_{i\in J_{e,h}}q_i\}
 \ge t_e^I/r.
\]
If $r=0$, the incoming gaps all vanish.

A different global law uses the same low success intervals and high
failure intervals $[0,q_i]$. Since $|O_e|\le s$, it gives
\[
 \E D_e\ge\min\{u_e,\max_{i\in O_e}q_i\}\ge t_e^O/s
\]
when $s>0$; empty outgoing sets have zero gap. Mixing these two laws and
independence with weights proportional to $r,s,\kappa$ proves the bound
$r+s+\kappa$, including zero parameters by omitting their components.

For the degree alternative, form the positive polynomial with terms
$a_eX_{i_e}\prod_{i\in O_e}X_i$ for one-low factors. Its degree is at most
$s+1$, its termwise gap is $\sum_e a_et_e^O$, and its deficiency under
\emph{every} common-marginal law is at most the original total deficiency.
Its hull gap is consequently at most the original hull gap. Thus the
outgoing sum is at most $B(s+1)\hgap$, as well as at most $s\hgap$.
The incoming sum is at most $r\hgap$ by the ownership law, and the other
terms' sum at most $\kappa\hgap$ by independence. Adding proves the
refinement. This argument needs only a polynomial-level degree bound,
not an unproved per-term guarantee from that bound. Orienting all edges
from variables to factors gives the frequency specialization.
\end{proof}

\begin{corollary}[Sharp ownership constant]\label{cor:ownership-sharp}
At points where every term has exactly one low coordinate, suppose each
high variable occurs in at most $r$ factors, with unrestricted low-variable
frequency. The sharp supremum of $\tgap/\hgap$ is $r$ for every integer
$r\ge1$.
\end{corollary}
The ownership proof gives the upper bound using only its $r$ rounds.
The lower bound for $r\ge2$ follows from the next construction with $L=r$;
for $r=1$ a single nonaffine product attains one. This parameter concerns
high variables at the stated point and is distinct from ordinary frequency.

\subsection{Variable radix and exact width}
\begin{theorem}[Variable-radix family]\label{thm:variable-radix}
Let $L\ge2$ and $b\ge L$ be integers, and put $m=b^L$. Index $m$ leaves
by length-$L$ base-$b$ words. Let $\mathcal P_j$ partition them by their
first $j$ digits. With separate anchors $a_1,\ldots,a_L$, define
\[
 f_{L,b}(a,z)=\sum_{j=1}^L\sum_{Q\in\mathcal P_j}
                       a_j\prod_{i\in Q}z_i,
 \qquad a_j=b^{-j},\quad z_i=1-m^{-1}.
\]
Then
\begin{equation}\label{eq:radix-gaps}
 \tgap=L,\qquad \hgap=1+(L-1)/b,
 \qquad \tw(G_{L,b})=L.
\end{equation}
\end{theorem}
\begin{proof}
A level-$j$ block has $m/b^j$ leaves, whose failure means sum to $b^{-j}$.
Its local lower value is zero and upper value $b^{-j}$; $b^j$ blocks give
one unit of termwise gap at each level. For an arbitrary binary law let
$R$ count failed leaves and $N_j$ count hit level-$j$ blocks. Then
$\E R=1$, $N_j\le\min(b^j,R)$, and total deficiency is
$\sum_j A_jN_j$. For a fixed distribution of $R$, a prescribed-mean
indicator has largest correlation with the nondecreasing function
$\min(b^j,R)$ by selecting its upper tail: exchanging selected probability
at a smaller count with unselected probability at a larger count proves
this assertion, with fractional selection at a tied count allowed. All
anchors can use these nested tails simultaneously. Represent the count by
its nonincreasing quantile $R(U)$ and select $A_j=\ind\{U\le b^{-j}\}$.
If exactly $l$ anchors are selected, then
\[
 \sum_{j=1}^l\min(b^j,R)\le R+\sum_{j=1}^{l-1}b^j.
\]
For $l=0$ use $0\le R$. Integrating gives
$\hgap\le1+\sum_{j=1}^{L-1}b^j b^{-j-1}=1+(L-1)/b$.

The probability of level $l$ is $w_l=(b-1)b^{-l-1}$ for $0\le l<L$
and $w_L=b^{-L}$. Two count profiles are
\[
 R_l^{(1)}=b^l\ind\{l\ge1\},\qquad
 R_l^{(2)}=b^{l-1}\ind\{l\ge2\}.
\]
Their means are $M_1=[(L-1)(b-1)+b]/b$ and
$M_2=[(L-2)(b-1)+b]/b^2$. Our parameter range ensures
$M_2\le1\le M_1$, so mixing the profiles with first-profile weight
$(1-M_2)/(M_1-M_2)$ gives mean count one. Each profile attains the
preceding pointwise inequality at every level, including $l=0,1$.
For an integer count $R$, fail the first $R$ leaves in digit-reversal
order. Their first $j$ digits are the reversals of the last $j$ digits of
$0,\ldots,R-1$, so they hit exactly $\min(b^j,R)$ level-$j$ blocks.
A uniform coordinatewise digit shift modulo $b$ preserves every hit count
and gives each leaf conditional failure probability $R/m$. The mixture
therefore has all the required marginals and attains the upper bound.
This argument works for composite radices as well as prime ones.

For treewidth, eliminate leaves first, filling their $L$ factor neighbors
into cliques. These neighbors form nested chains. Next eliminate factor
vertices from level $L$ down to level one. A level-$j$ factor's remaining
neighbors lie among its $j-1$ ancestor factors and anchors of levels
$j,\ldots,L$, at most $L$ vertices. Induction verifies this description:
eliminating a descendant joins its ancestor chain and deeper anchors,
never two unrelated factors. The final $L$ anchors also have elimination
width at most $L$. For the reverse bound choose a level-$(L-1)$ block
with $b$ leaves and contract each deepest bilinear factor into its leaf.
The resulting minor contains $K_{L,b}$, with the $L-1$ ancestor factors
and $a_L$ opposite those leaves. For $b\ge L$, this graph has a clique
minor of order $L+1$: pair $L-1$ vertices on opposite sides into branch
sets, then keep one unused vertex from each side. All branch sets are
connected and pairwise adjacent. Minor monotonicity proves the lower bound.
\end{proof}

The restriction $b\ge L$ is part of the exact formula just proved. In
particular it cannot be replaced silently by $b=2$ for arbitrarily many
levels; Theorem~\ref{thm:dyadic-exact} uses a different cutoff. More
generally, for integers $b\ge2,L\ge2$, the same radix proof gives
$\hgap=s+(L-s)b^{-s}$ whenever $1\le s<L$ and
\[
 M_{s+1}\le1\le M_s,\qquad
 M_q=\frac{(L-q)(b-1)+b}{b^q}.
\]
Indeed, cap the first $l-s$ summands and bound the last $s$ by $R$, giving
$\sum_{j\le l}\min(b^j,R)\le sR+\sum_{j\le l-s}b^j$.
Its integrated constant is $(L-s)b^{-s}$. The profiles
$R_l^{(q)}=b^{l-q+1}\ind\{l\ge q\}$ for $q=s,s+1$ have means $M_q$
and attain equality; mix them to mean one and use digit reversal as above.

\begin{corollary}[Sharp leading structural growth]\label{cor:structural-growth}
Let $W_k,D_k,P_k$ be the supremal unit-cube positive-monomial gap ratios
under, respectively, incidence treewidth, degeneracy, or orientation
parameter at most $k$. For every integer $k\ge2$,
\[
 W_k\ge k,\qquad D_k\ge k+1,\qquad P_k\ge k+1,
\]
whereas
\[
 W_k\le D_k\le P_k\le k+\min\{k,B(k+1)\}+\kappa.
\]
In particular $W_k\sim D_k\sim P_k\sim k$ as $k\to\infty$.
\end{corollary}
\begin{proof}
Treewidth at most $k$ implies degeneracy at most $k$: in a reduced tree
decomposition take a vertex exclusive to a leaf bag and delete it,
then continue on induced subgraphs. A degeneracy order oriented toward
later vertices has outdegrees at most $k$. This proves the upper chain.
Use $L=k$ in Theorem~\ref{thm:variable-radix} and let $b\to\infty$ for
$W_k\ge k$. For the stronger other bounds use $L=k+1$: orient each
deepest factor toward both its leaf and anchor, all other leaf incidences
toward their factor, and the remaining factor incidences toward their
anchor. Outdegrees are at most $k$, since deepest factors have outdegree
two. For degeneracy delete deepest factors of degree two, their isolated
anchor, then leaves of degree $L-1=k$, leaving stars. Letting $b$ grow
proves both lower bounds. This family still has treewidth $k+1$.
For example $L=3,b=5$ gives $15/7>2$ for degeneracy and orientation
parameter two, with treewidth three. Finally $B(k+1)=o(k)$.
\end{proof}

For ordinary frequency, the example in Proposition~\ref{prop:feedback-sharp}
below gives a lower bound $2-1/\nu$ for every $\nu\ge2$.
The dyadic family has maximum frequency $2^L$ (at its deepest anchor), so
choosing the largest power of two at most $\nu$ also gives the lower
asymptotic $(1-o(1))\log\nu/\log\log\nu$. Frequency two has the sharper
exact value $3/2$, proved in Section~\ref{sec:frequency}.

\subsection{A lower bound on the normalized means}
Let $C_\delta$ be the supremal positive-monomial gap ratio on unit cubes
at points whose coordinates are all at least $\delta$, with dimension
and degree unrestricted. This restricts the \emph{evaluation point}; the
envelopes remain those on the full unit cube.

\begin{theorem}[Marginal-floor law]\label{thm:marginal-floor}
For $0<\delta\le1/2$, set
\[
 M=\max\{2,\log(1/\delta)\},\quad \tau=\delta/M^2,\quad
 \Lambda=\log\frac{1+\tau}{\tau},\quad
 I_\delta=\int_0^1\left(1-e^{-1/[\Lambda(z+M^{-2})]}\right)\,dz.
\]
Then $I_\delta>0$ and $C_\delta\le I_\delta^{-1}+\kappa$.
For $1/2<\delta<1$, $C_\delta\le C_{1/2}$. Moreover
\[
 C_\delta\sim\frac{\log(1/\delta)}{\log\log(1/\delta)}
 \quad(\delta\downarrow0).
\]
The same leading asymptotic holds with every coordinate in
$[\delta,1-\delta]$. The upper bound transfers to nonnegative boxes when
all nonfixed normalized means $(x_i-\ell_i)/(u_i-\ell_i)$ are at least
$\delta$.
\end{theorem}
\begin{proof}
The density $h(t)=1/[\Lambda(t+\tau)]$ integrates to one on $(0,1)$.
Use low-coordinate success intervals $[0,x_i]$. For a high coordinate
with failure mean $q<1/2$ define
\[
 q_0(t)=\min\{1,qh(t)\},\quad m_q=\int_0^1q_0(t)\,dt,\quad
 q_1(t)=q_0(t)+\frac{q-m_q}{1-m_q}(1-q_0(t)).
\]
Here $0\le m_q\le q<1$, so $q_1\in[q_0,1]$ and its integral is exactly
$q$, including $q=0$. Conditional on the common uniform variable $U=t$,
sample high failures independently with probabilities $q_1(t)$.
For a set of high coordinates with failure sum $S$, the probability of
some failure is at least $1-e^{-Sh(t)}$. If any $qh(t)\ge1$, a failure
is certain; otherwise $q_1(t)\ge qh(t)$ and the usual product bound proves
it. The clipping case is needed for this conclusion.

A one-low term with anchor mean $u\ge\delta$ consequently has
\[
 \E D_e\ge\int_0^u(1-e^{-S/[\Lambda(t+\tau)]})\,dt.
\]
For $S>0$, substitute $t=uz$, $y=S/u$, and note
\[
 \frac{1-e^{-ay}}{\min(1,y)}\ge1-e^{-a}\quad(a\ge0,y>0).
\]
Concavity proves this for $y\le1$, monotonicity for $y\ge1$.
Since $\tau/u\le M^{-2}$, division by $t_e=u\min(1,y)$ now gives
$\E D_e/t_e\ge I_\delta$. Zero $S$ requires no division. Mix this law
and independence with weights proportional to $I_\delta^{-1},\kappa$;
Lemma~\ref{lem:easy-terms} handles all other factors.

As $\delta\downarrow0$, $\Lambda=M+2\log M+o(1)$ and
$\Lambda/M^2\to0$. The inequalities $1-e^{-a}\le\min(1,a)$ and
$1-e^{-a}\ge a-a^2/2$ give
\[
 I_\delta\le\frac{1+\log\Lambda}{\Lambda},\qquad
 I_\delta\ge\frac{\log\Lambda}{\Lambda(1+\Lambda/M^2)}
                   -\frac{\Lambda-1}{2\Lambda^2}.
\]
For the second bound integrate only over $[1/\Lambda,1]$, using
$z+M^{-2}\le z(1+\Lambda/M^2)$. Hence
$I_\delta\sim\log\Lambda/\Lambda$, proving the upper asymptotic.
The dyadic family with $L=\lfloor\log_2(1/\delta)\rfloor$ has all means
in $[\delta,1-\delta]$ and ratio asymptotic to $L/\log_2 L$, which gives
the matching lower bound. Monotonicity gives the finite large-$\delta$
claim. Finally positive affine expansion on a nonnegative box preserves
the lower restriction on normalized means, so
Proposition~\ref{prop:box-transfer} gives its stated extension.
\end{proof}

\begin{theorem}[Uniform joint asymptotic]\label{thm:joint-growth}
Let $C(n,d,k,\delta)$ be the supremal unit-cube positive-monomial ratio
in dimension at most $n$, degree at most $d$, incidence treewidth at most
$k$, and with every mean at least $\delta$. Assume integers $n,d\ge2$,
$k\ge1$, and $0<\delta\le1/2$. Put
$q=\min\{n,d,1/\delta\}\ge e^e$,
$\phi(q)=\log q/\log\log q$, and $h=\min\{k,\phi(q)\}$.
Along every joint parameter sequence with $h\to\infty$,
\[
 C(n,d,k,\delta)\sim h.
\]
The same statement holds for the two-sided strip and with degeneracy or
the orientation parameter replacing treewidth.
\end{theorem}
\begin{proof}
The separate dimension, degree, marginal-floor, and incidence estimates
give $C\le\min\{(1+o(1))\phi(q),k+o(k)\}$. Their relative errors vanish
uniformly along the stated sequences, since both $q$ and $k$ then diverge.
For a simultaneous lower construction take
\[
 L=\lfloor h\rfloor,\qquad b=\lfloor(q/2)^{1/L}\rfloor,\qquad m=b^L.
\]
Then $b\ge(1-o(1))\log q$, $L/b\le(1+o(1))/\log\log q\to0$,
and eventually $b\ge L\ge2$. The radix family has dimension
$m+L\le q/2+\phi(q)\le q\le n$, degree $m/b+1\le q\le d$, and
width $L\le k$. Its smallest anchor mean and leaf failure mean are
$1/m\ge2/q\ge\delta$, and all other means also lie in the two-sided
strip. Thus its ratio $L/[1+(L-1)/b]\sim h$. Its degeneracy and
orientation parameter are at most its treewidth, proving the alternatives.
This is not an asymptotic assertion for a fixed small value of $k$.
\end{proof}
